% =====================================================================
\chapter*{Instrucciones}
\addcontentsline{toc}{chapter}{Instrucciones}
\markboth{Instrucciones}{}
% =====================================================================

Este examen evalúa la fase~IV del curso: preentrenamiento a escala (módulo~16), hardware y sistemas (módulo~17), arquitecturas modernas (módulo~18) e inferencia (módulo~19). Casi todas las preguntas son estimaciones de orden de magnitud; se valora más el razonamiento, con unidades en cada paso, que el último decimal.

\section*{Condiciones}
\begin{center}
\small
\begin{tabularx}{\textwidth}{@{}lX@{}}
\toprule
Duración & 3 horas. \\
Material permitido & Calculadora científica no programable y una hoja A4 de fórmulas escrita a mano por ambas caras. No se permiten apuntes, libros ni dispositivos conectados. \\
Puntuación & 100 puntos: Parte~A, 40 puntos (20 cuestiones de 2 puntos); Parte~B, 60 puntos (6 problemas de 10 puntos). \\
Aprobado & 70 puntos (70\,\%) o más. \\
Módulos & M16, M17, M18 y M19 (véase el reparto por módulo en el \cref{ap:autoevaluacion}). \\
Proyecto de fase & Se evalúa aparte sobre 100 puntos (\cref{cap:proyecto}) y se entrega en el plazo que fije el profesor. \\
\bottomrule
\end{tabularx}
\end{center}

\section*{Normas}
\begin{enumerate}[label=\arabic*.]
  \item En la Parte~A hay una sola opción correcta por cuestión. Las respuestas erróneas no restan.
  \item En la Parte~B deben escribirse las unidades en cada paso. Los resultados finales se dan con tres cifras significativas; los intermedios, con cuatro.
  \item Salvo que el enunciado diga otra cosa, se usan los datos de referencia del \cref{tab:datos} y las fórmulas que siguen. Las cifras de hardware son aproximadas y proceden de hojas de especificaciones públicas; las de modelos reales, de sus artículos o fichas técnicas públicas. Los modelos «Cernícalo-MoE», «Halcón-35» y «Alcotán» son inventados.
  \item Se toma $1~\text{GB}=10^{9}$~bytes y $1~\text{KiB}=1024$~bytes, $1~\text{GiB}=2^{30}$~bytes.
\end{enumerate}

\section*{Datos de referencia}
\begin{table}[h]
\centering\small
\caption{Datos de referencia del examen (valores aproximados).}\label{tab:datos}
\begin{tabularx}{\textwidth}{@{}lX@{}}
\toprule
Dato & Valor \\
\midrule
NVIDIA H100 SXM & 989~TFLOP/s en bf16 densos; 80~GB de HBM3; 3{,}35~TB/s; hasta 700~W \\
NVIDIA A100 SXM 80~GB & 312~TFLOP/s en bf16 densos; 80~GB; 2~TB/s; 400~W \\
NVLink (H100, dentro de un nodo de 8 GPU) & $\approx 900$~GB/s por GPU \\
InfiniBand entre nodos & $\approx 50$~GB/s por GPU (400~Gb/s) \\
Bytes por parámetro con Adam y precisión mixta & 16 (2 de pesos, 2 de gradientes y 12 del optimizador) \\
Ley de Chinchilla (enfoque~3) & $L(N,D)=E+A/N^{\alpha}+B/D^{\beta}$, con $E=1{,}69$, $A=406{,}4$, $B=410{,}7$, $\alpha=0{,}34$, $\beta=0{,}28$ \\
\bottomrule
\end{tabularx}
\end{table}

\section*{Fórmulas}
\[
C\approx 6ND,\qquad
t=\frac{C}{n_{\mathrm{d}}\,\Ppico\,\MFU},\qquad
\ckv=2\,L\,\nkv\,\dhd\,s\ \ \text{bytes por token},\qquad
\Icrit=\frac{\Ppico}{\Bmem},
\]
donde $n_{\mathrm{d}}$ es el número de dispositivos, $L$ el número de capas, $\nkv$ el de cabezas de clave y valor, $\dhd$ la dimensión de cada cabeza y $s$ los bytes por número (2 en bf16).
